\section{Chapter Summary}
\label{chap:sota:summary}

This chapter reviewed the state of the art on the technologies, standards and experimental platforms related to the exposure and testing of \acs{5G} capabilities and network-aware applications. Section~\ref{sota:5gecosystem} began with an overview of the current state of the \acs{5G} ecosystem. It presents its three traditional service classes, \acs{eMBB}, \acs{URLLC} and \acs{mMTC}, along with network slicing as the mechanism to accommodate their heterogeneous \acs{QoS} requirements over a shared infrastructure, supported by virtualization and softwarization technologies such as \acs{NFV} and \acs{SDN}.

Section~\ref{sota:slicing} further detailed network slicing, covering the use of the \acs{GST} and \acsp{NEST} to describe network slice requirements and the phases of the \acs{NSI} lifecycle. It then described the hierarchy of management functions, the \acs{CSMF}, \acs{NSMF} and \acs{NSSMF}, capable of orchestrating \acsp{NSI} and their constituent \acsp{NSSI}. It also presented the notion of slice managers, software systems that automate the orchestration and lifecycle management of network slices, with particular focus on the Katana and the \acs{ITAv} Slice Managers. Section~\ref{sota:nfv} then focused on the \acs{NFV-MANO} architectural framework and its functional blocks.

Afterwards, the \acs{5G} \acs{SBA} and its constituent network functions were presented in Section~\ref{sota:sba}, highlighting the analytics and exposure capabilities of the \acs{5G} System, as specified by \acs{3GPP}. The exposure capabilities comprise the \acs{NEF} and the Application Enabler Layer \acsp{API}. In order to assess the availability of these interfaces in practice, the existing \acs{NEF} implementations were reviewed. This evaluation concluded that \acs{5GC} available open source solutions typically lack a \acs{NEF} or provide only a partial implementation. Among the analyzed solutions, \acs{NEF}Sim stood out by allowing network applications to be implemented and validated without requiring access to a full \acs{NEF} instance.

Section~\ref{sota:networkapplications} discussed the different definitions of network applications, converging on application services that provide functionalities to vertical industries by interacting with the \acs{5G} control plane through the aforementioned exposure interfaces. It also outlined use cases of such applications reported in the literature, in the \acs{CAM}, \acs{TL}, media streaming and \acs{PPDR} sectors. Section~\ref{sota:telcoinitiatives} then addressed the telecom standardization initiatives relevant to these applications. The CAMARA Project and the \acs{GSMA} Open Gateway define simplified network \acsp{API} for developers without telecom expertise, while \acs{CAPIF} provides a single framework for the secure publication, discovery and invocation of northbound \acsp{API}. The \acs{TMF} Open \acsp{API} and the \acs{SID} model were introduced as well, focusing on the Product, Service and Resource domains.

Section~\ref{sota:bss-oss-mano} described the role of \acsp{BSS} and \acsp{OSS} in the management of telecom operations, with the former handling customer-facing interactions and the latter the underlying network execution. It also outlined the benefits of integrating these systems with \acs{NFV-MANO} solutions through standardized interfaces.

Regarding experimentation environments, Section~\ref{sec:testbeds:main} analyzed experimentation environments from diverse works throughout the literature. These environments follow a common architecture, combining open source or commercial \acs{5G} cores and \acs{RAN} implementations with virtualization and orchestration platforms. In terms of network exposure, the analysis showed varying levels of adoption of \acs{NEF}, CAMARA, \acs{CAPIF} and \acs{TMF} \acsp{API}, with none of the analyzed solutions combining all four capabilities.

To understand how network-aware applications and their supporting components are evaluated in practice, the corresponding testing and validation procedures were reviewed in Section~\ref{sec:testingandvalidation}. Validation of network applications was mostly focused on \acs{QoS} management scenarios, assessing throughput under the applied \acs{QoS} profiles. The testing pipelines found in the literature present different levels of automation, ranging from partially manual procedures to multistage pipelines orchestrated through \acsp{OSS}. The section concluded with a consolidated view of the metrics and \acsp{KPI} collected and evaluated across the analyzed works.
